\documentclass[10pt,a4paper]{amsart}
\usepackage[margin=19mm]{geometry}
\usepackage{amsmath,amssymb,mathtools}
\usepackage[scr=rsfs,cal=euler]{mathalfa}
\usepackage{xfrac}
\usepackage[unicode,hidelinks,bookmarks=false]{hyperref}
\newcommand{\ie}{i.e.\ }
\newcommand{\cf}{cf.\ }
\newcommand{\resp}{respectively\ }
\newcommand{\Subsection}{\subsection}
\providecommand{\nobreakdash}{\nobreak\hbox{-}}
\setlength{\emergencystretch}{2em}
\multlinegap0pt
\usepackage{bm}
\usepackage{bbm}
\usepackage{dsfont}
\usepackage{xspace}
\usepackage{cancel}
\usepackage{mathtools}
\usepackage{amsthm}
\makeatletter
\@ifundefined{theorem}{\newtheorem{theorem}{Theorem}}{}
\@ifundefined{lemma}{\newtheorem{lemma}[theorem]{Lemma}}{}
\makeatother
\def\N{\mathop{\mathbb N\kern 0pt}\nolimits}
\def\R{\mathop{\mathbb R\kern 0pt}\nolimits}
\def\cQ{\mathcal Q}
\def\ds{\displaystyle}
\def\f{\frac}
\def\ls{\lesssim}
\def\na{\nabla}
\def\p{\partial}
\def\ve{\varepsilon}
\title[Global classical solutions to 3D irrotational compressible E]{Global classical solutions to 3D irrotational compressible Euler equations of Chaplygin gases with weakly decaying initial data}
\author{Mu Gao, Huicheng Yin}
\date{Source: arXiv:2608.13166v1 (CC BY 4.0)}
\begin{document}
\maketitle

\noindent\emph{Source and licence.} Global classical solutions to 3D irrotational compressible Euler equations of Chaplygin gases with weakly decaying initial data. Version arXiv:2608.13166v1 (CC BY 4.0), distributed under the Creative Commons Attribution 4.0 International licence, as recorded in the arXiv licence metadata for this exact version and shown on the abstract page https://arxiv.org/abs/2608.13166v1. Full rights notice: licenses/arxiv-2608.13166-CC-BY-4.0.txt.\par\smallskip

\noindent\textit{Editorial scope.} Counted unit: the Lemma whose source label is \texttt{lem:Chaplygin-energy-estimate} in the article (source0.tex lines 4404--4415, source label \texttt{lem:Chaplygin-energy-estimate}), delivered together with the article's own complete original proof of it (source0.tex lines 4417--4815, one \texttt{proof} environment of 399 lines). No mathematics is written, repaired or re-derived: statement and proof are the article's own text line for line. Required context retained uncounted: YHC-4 (lines 447--451); eq:Chaplygin-Cauchy (lines 460--469); eq:CVP-bootstrap-parameter-range (lines 3538--3545); eq:CVP-bootstrap-energy-ordinary (lines 3549--3556); eq:CVP-bootstrap-pointwise-phi (lines 3568--3579); eq:CVP-bootstrap-Strichartz-Gamma (lines 3592--3608); eq:Chaplygin-energy-reformulation (lines 4381--4386); eq:Chaplygin-energy-symmetry (lines 4396--4402). Every definition, display and label of those lines is retained unchanged. Omitted from this package and not redistributed: 1--446, 452--459, 470--3537, 3546--3548, 3557--3567, 3580--3591, 3609--4380, 4387--4395, 4403--4403, 4416--4416, 4816--8496. Nothing inside a retained excerpt is omitted or altered: every retained body is delivered exactly as the article's lines are written, so no omission receipt of any kind accompanies this package. Citations: the retained excerpts carry no citation command at all, so no bibliography entry of the article is transcribed and no reference is redistributed. Reference adaptation: none was needed and none was made; every reference inside the delivered unit resolves inside it, so no unresolved reference or citation is produced. Printed slips kept literal: no printed word, symbol, hypothesis or number of the counted statement or of its proof was altered. Layout and class: the article's own class is \texttt{article} with packages including amsmath, amssymb, amsthm, mathrsfs, simplewick, times, color, hep, leftidx, graphicx, float, enumerate, titletoc, epstopdf; the wrapper is the lane's pinned \texttt{amsart} editorial wrapper at 10pt on A4, which restates the theorem-like environments the retained text needs and adds the article's own macro definitions that the retained text needs (\texttt{\textbackslash N}, \texttt{\textbackslash R}, \texttt{\textbackslash cQ}, \texttt{\textbackslash ds}, \texttt{\textbackslash f}, \texttt{\textbackslash ls}, \texttt{\textbackslash na}, \texttt{\textbackslash p}, \texttt{\textbackslash ve}), with extra wrapper packages (bm, bbm, dsfont, xspace, cancel, mathtools). The delivered numbering is the unit's own. Assets: no retained span contains a figure environment, a graphics-inclusion command, a PDF-page inclusion, a table or a TikZ picture; no image file is copied into this package, the package declares no asset dependency and no image file is redistributed with it. Licence: version arXiv:2608.13166v1 of this article is distributed under the Creative Commons Attribution 4.0 International licence, as recorded in the arXiv licence metadata for this exact version and shown on the abstract page https://arxiv.org/abs/2608.13166v1. A full rights notice is delivered at licenses/arxiv-2608.13166-CC-BY-4.0.txt. Characters: every retained excerpt is pure ASCII, so the delivered engine, XeLaTeX with its default Latin Modern text font, reports no missing character.
\begin{equation}\label{YHC-4}
\begin{split}
\|\phi_0\|_{H^{N+1}(\mathbb{R}^3)}+\|\phi_1\|_{H^{N}(\mathbb{R}^3)} + \sum_{|a|\le 13} \|\langle x\rangle^{1+\mu} \nabla^a(\na\phi_0,\phi_1)\|_{L^2(\mathbb{R}^3)} \le C_0,
\end{split}
\end{equation}

\begin{equation}\label{eq:Chaplygin-Cauchy}
\left\{
\begin{aligned}
&\Box\phi=2\partial_t\phi\,\Delta\phi-
2\sum_{k=1}^3\partial_k\phi\,\partial^2_{tk}\phi+
|\nabla\phi|^2\Delta\phi-\sum_{k,j=1}^3\partial_k\phi\,\partial_j\phi\,\partial^2_{kj}\phi,\\
&(\phi, \partial_t\phi)(0,x)=\ve(\phi_0(x),\phi_1(x)).
\end{aligned}
\right.
\end{equation}

		\begin{equation}
		\label{eq:CVP-bootstrap-parameter-range}
		N \ge 15,
		\qquad
		0<\mu<\frac12,
		\qquad
		0\le\eta\le\frac{\mu}{100}.
	\end{equation}

	\begin{equation}
		\label{eq:CVP-bootstrap-energy-ordinary}
		\sup_{0\le t\le T}
		(1+t)^{-\eta}
		\|\p\phi(t)\|_{H^N(\R^3)}
		\le
		\ve_1,
	\end{equation}

	\begin{equation}\label{eq:CVP-bootstrap-pointwise-phi}
		\big\|
		AB^{\f\mu4}\phi
		\big\|_{L^\infty([0,T]\times\R^3)}
		+
		\sum_{|a|\le 7}
		\big\|
		AB^{\f\mu4}\nabla^a\p\phi
		\big\|_{L^\infty([0,T]\times\R^3)}
		\le
		\ve_1,
	\end{equation}

	\begin{equation}
		\label{eq:CVP-bootstrap-Strichartz-Gamma}
		\sum_{\Gamma}
		\big\|
		A^{\frac{\mu}{3}}
		\Gamma\phi
		\big\|_{L^2([0,T];L^\infty(\mathbb R^3))}
		+
		\sum_{\Gamma}
		\sum_{|a|\le 6}
		\big\|
		A^{\frac{\mu}{3}}
		\partial\nabla^a\Gamma\phi
		\big\|_{L^2([0,T];L^\infty(\mathbb R^3))}
		\le
		\varepsilon_1.
	\end{equation}

	\begin{equation}\label{eq:Chaplygin-energy-reformulation}
		\Box\phi
		=\sum_{\alpha,\beta=0}^3
		\cQ^{\alpha\beta}(\partial\phi)
		\partial_{\alpha\beta}^2\phi,
	\end{equation}

	\begin{equation}\label{eq:Chaplygin-energy-symmetry}
		\cQ^{\alpha\beta}
		=
		\cQ^{\beta\alpha},
		\qquad
		0\le\alpha,\beta\le3.
	\end{equation}

	\begin{lemma}[{\bf Energy estimate}]\label{lem:Chaplygin-energy-estimate}
		Let \(\phi\) be a smooth solution to
		\eqref{eq:Chaplygin-Cauchy} on \([0,T]\times\mathbb R^3\).
		Assume that the initial data satisfy \eqref{YHC-4} and
		\eqref{eq:CVP-bootstrap-parameter-range}-\eqref{eq:CVP-bootstrap-Strichartz-Gamma}
		hold. Suppose that \(\ve_1>0\) is small. Then
		\begin{equation}\label{eq:Chaplygin-energy-growth-epsilon1}
			\sup_{0\le t\le T}(1+t)^{-\eta}\|\partial\phi(t)\|_{H^N(\mathbb R^3)}
			\ls
			\ve +\ve_1^2.
		\end{equation}
	\end{lemma}

	\begin{proof}
		Acting $\partial_x^a$ with $a\in\N_0^3$ on two sides of
		\eqref{eq:Chaplygin-energy-reformulation} yields
		\begin{equation}\label{eq:Chaplygin-commuted-equation}
			\begin{aligned}
				\Box\partial_x^a\phi
				&=
				I_1^a+I_2^a,
				\\
				I_1^a
				=
				\sum_{\alpha,\beta=0}^3
				\cQ^{\alpha\beta}
				\partial_x^a\partial_{\alpha\beta}^2\phi,\quad
				&I_2^a
				=
				\sum_{\alpha,\beta=0}^3
				\sum_{\substack{b+c=a,\\ |b|\le |a|-1}}
				C_{abc}
				\partial_x^c\cQ^{\alpha\beta}
				\partial_x^b\partial_{\alpha\beta}^2\phi,
			\end{aligned}
		\end{equation}
		where $I_2^a$ vanishes for $|a|=0$.
		
		Multiplying \eqref{eq:Chaplygin-commuted-equation} by
		$\partial_t\partial_x^a\phi$ and integrating the resulting equality over
		$[0,t]\times\mathbb R^3$, one has
		\begin{equation}\label{eq:Chaplygin-basic-energy}
			\begin{aligned}
				\|\partial_x^a\partial\phi(t)\|_{L^2(\mathbb R^3)}^2
				&\ls
				\|\partial_x^a\partial\phi(0)\|_{L^2(\mathbb R^3)}^2
				+
				\big|
				\int_0^t\int_{\mathbb R^3}
				\partial_t\partial_x^a\phi\,
				I_1^a\,dx\,d\tau
				\big|
				\\
				&\quad+
				\int_0^t
				\|\partial_t\partial_x^a\phi(\tau)\|_{L^2(\mathbb R^3)}
				\|I_2^a(\tau)\|_{L^2(\mathbb R^3)}
				d\tau.
			\end{aligned}
		\end{equation}
		At first, we deal with the last term in the first two lines of
		\eqref{eq:Chaplygin-basic-energy}. Note that
		\begin{equation}\label{eq:Chaplygin-principal-identity}
			\begin{aligned}
				\cQ^{\alpha\beta}
				\partial_t\partial_x^a\phi\,
				\partial_{\alpha\beta}^2\partial_x^a\phi
				&=
				\partial_\alpha
				\big(
				\cQ^{\alpha\beta}
				\partial_t\partial_x^a\phi\,
				\partial_\beta\partial_x^a\phi
				\big)
				-
				\partial_\alpha\cQ^{\alpha\beta}
				\partial_t\partial_x^a\phi\,
				\partial_\beta\partial_x^a\phi
				\\
				&\quad-
				\frac12\partial_t
				\big(
				\cQ^{\alpha\beta}
				\partial_\alpha\partial_x^a\phi\,
				\partial_\beta\partial_x^a\phi
				\big)
				+
				\frac12\partial_t\cQ^{\alpha\beta}
				\partial_\alpha\partial_x^a\phi\,
				\partial_\beta\partial_x^a\phi,
			\end{aligned}
		\end{equation}
		where the Einstein summation convention for
		$\alpha,\beta=0,1,2,3$ and the symmetric condition
		\eqref{eq:Chaplygin-energy-symmetry} are used in
		\eqref{eq:Chaplygin-principal-identity}.
		
		Due to
		\begin{equation}\label{eq:Chaplygin-Q-pointwise}
			|\cQ^{\alpha\beta}|
			\ls
			|\partial\phi|
			+
			|\partial\phi|^2,
		\end{equation}
		it follows from \eqref{eq:CVP-bootstrap-pointwise-phi} that
		\begin{equation}\label{eq:Chaplygin-Q-Linfty}
			\begin{aligned}
				\|\cQ^{\alpha\beta}(t)\|_{L^\infty}
				&\ls
				\|\partial\phi(t)\|_{L^\infty}
				+
				\|\partial\phi(t)\|_{L^\infty}^2
				\\
				&\ls
				\|\partial\phi(t)\|_{W^{2,\infty}}
				+
				\|\partial\phi(t)\|_{W^{2,\infty}}^2
				\\
				&\ls
				(1+t)^{-1}\varepsilon_1
				+
				(1+t)^{-2}\varepsilon_1^2
				\\
				&\ls
				\varepsilon_1.
			\end{aligned}
		\end{equation}
		
		Direct computation yields
		\begin{equation}\label{eq:Chaplygin-spatial-Q-derivative}
			\begin{aligned}
				\|\nabla\cQ^{\alpha\beta}(t)\|_{L^\infty}
				&\ls
				\|\nabla\partial\phi(t)\|_{L^\infty}
				+
				\|\partial\phi(t)\|_{L^\infty}
				\|\nabla\partial\phi(t)\|_{L^\infty}
				\\
				&\ls
				\|\partial\phi(t)\|_{B_{\infty,1}^{1+\delta}}
				+
				\|\partial\phi(t)\|_{B_{\infty,1}^{1+\delta}}^2.
			\end{aligned}
		\end{equation}
		Next we estimate $\partial_t\cQ^{\alpha\beta}$. Due to
		$\partial_t^2\phi
			=\Delta\phi	+2\ds\sum_{j=1}^3\cQ^{0j}\partial_{tj}^2\phi
			+\ds\sum_{k,j=1}^3\cQ^{kj}\partial_{kj}^2\phi$,
		by \eqref{eq:Chaplygin-Q-Linfty}, one has
		\begin{equation}\label{eq:Chaplygin-ttphi-Linfty}
			\begin{aligned}
				\|\partial_t^2\phi(t)\|_{L^\infty}
				&\ls
				\|\nabla\partial\phi(t)\|_{L^\infty}
				+
				\|\cQ(t)\|_{L^\infty}
				\|\nabla\partial\phi(t)\|_{L^\infty}
				\\
				&\ls
				\|\partial\phi(t)\|_{B_{\infty,1}^{1+\delta}}
				+
				\|\partial\phi(t)\|_{B_{\infty,1}^{1+\delta}}^2.
			\end{aligned}
		\end{equation}
		Consequently,
		\begin{equation}\label{eq:Chaplygin-time-Q-derivative}
			\begin{aligned}
				\|\partial_t\cQ^{\alpha\beta}(t)\|_{L^\infty}
				&\ls
				\|\partial^2\phi(t)\|_{L^\infty}
				+\|\partial\phi(t)\|_{L^\infty}
				\|\partial^2\phi(t)\|_{L^\infty}\\
				&\ls\|\partial\phi(t)\|_{B_{\infty,1}^{1+\delta}}
				+\|\partial\phi(t)\|_{B_{\infty,1}^{1+\delta}}^2.
			\end{aligned}
		\end{equation}
		Combining
		\eqref{eq:Chaplygin-spatial-Q-derivative} and
		\eqref{eq:Chaplygin-time-Q-derivative}, one has
		\begin{equation}\label{eq:Chaplygin-full-Q-derivative}
			\|\partial\cQ^{\alpha\beta}(t)\|_{L^\infty}
			\ls
			\|\partial\phi(t)\|_{B_{\infty,1}^{1+\delta}}
			+
			\|\partial\phi(t)\|_{B_{\infty,1}^{1+\delta}}^2.
		\end{equation}
		
		It follows from \eqref{eq:Chaplygin-Q-Linfty},
		\eqref{eq:Chaplygin-full-Q-derivative} and integration over
		$[0,t]\times\mathbb R^3$ for
		\eqref{eq:Chaplygin-principal-identity} that
		\begin{equation}\label{eq:Chaplygin-I1-estimate}
			\begin{aligned}
				\big|
				\int_0^t\int_{\mathbb R^3}
				\partial_t\partial_x^a\phi\,
				I_1^a\,dx\,d\tau
				\big|
				&\ls
				\varepsilon_1
				\|\partial_x^a\partial\phi(0)\|_{L^2}^2
				+
				\varepsilon_1
				\|\partial_x^a\partial\phi(t)\|_{L^2}^2
				\\
				&\quad+
				\int_0^t
				\big(
				\|\partial\phi(\tau)\|_{B_{\infty,1}^{1+\delta}}
				+
				\|\partial\phi(\tau)\|_{B_{\infty,1}^{1+\delta}}^2
				\big)
				\|\partial_x^a\partial\phi(\tau)\|_{L^2}^2
				d\tau.
			\end{aligned}
		\end{equation}
		
		
		Next, we treat $\|I_2^a\|_{L^2(\mathbb R^3)}$ with $|a|\ge1$ in
		\eqref{eq:Chaplygin-commuted-equation}. It is easy to find that
		\begin{equation}\label{eq:Chaplygin-I2-frequency}
			\begin{aligned}
				\|I_2^a\|_{L^2}
				&\ls
				\big\|
				\|P_kI_2^a\|_{L^2}
				\big\|_{\ell_k^2}\\
				&\ls
				\sum_{\alpha,\beta=0}^3
				\sum_{\substack{b+c=a,\\ |b|\le |a|-1}}
				\big\|
				\sum_{(k_1,k_2)\in\mathcal X_k}
				\big\|
				P_k
				\big(
				P_{k_1}\partial_{\alpha\beta}^2\partial_x^b\phi
				P_{k_2}\partial_x^c\cQ^{\alpha\beta}
				\big)
				\big\|_{L^2}
				\big\|_{\ell_k^2}.
			\end{aligned}
		\end{equation}
		Note that $\cQ^{00}=0$ and $|c|\ge1$ in \eqref{eq:Chaplygin-I2-frequency}. When
		$\max\{k_1,k_2\}=k_1$, one has that by Bernstein's inequality
		\begin{equation}\label{eq:Chaplygin-I2-first-frequency-case}
			\begin{aligned}
				&
				\big\|
				P_k
				\big(
				P_{k_1}\partial_{\alpha\beta}^2\partial_x^b\phi
				P_{k_2}\partial_x^c\cQ^{\alpha\beta}
				\big)
				\big\|_{L^2}
				\\
				&\ls
				\|P_{k_1}\partial\partial_x\partial_x^b\phi\|_{L^2}
				\|P_{k_2}\partial_x^c\cQ^{\alpha\beta}\|_{L^\infty}
				\\
				&\ls
				2^{k_1(|b|+1)}
				2^{k_2(|c|-1)}
				\|P_{k_1}\partial\phi\|_{L^2}
				\|P_{k_2}\partial_x\cQ^{\alpha\beta}\|_{L^\infty}
				\\
				&\ls
				2^{k_1|a|}
				\|P_{k_1}\partial\phi\|_{L^2}
				\|P_{k_2}\partial_x\cQ^{\alpha\beta}\|_{L^\infty}.
			\end{aligned}
		\end{equation}
		Therefore,
		\begin{equation}\label{eq:Chaplygin-I2-first-frequency-sum}
			\begin{aligned}
				&
				\big\|
				\sum_{\substack{(k_1,k_2)\in\mathcal X_k,\\
				\max\{k_1,k_2\}=k_1}}
				\big\|
				P_k
				\big(
				P_{k_1}\partial_{\alpha\beta}^2\partial_x^b\phi
				P_{k_2}\partial_x^c\cQ^{\alpha\beta}
				\big)
				\big\|_{L^2}
				\big\|_{\ell_k^2}
				\ls
				\|\partial\phi\|_{H^{|a|}}
				\|\cQ^{\alpha\beta}\|_{B_{\infty,1}^{1+\delta}}.
			\end{aligned}
		\end{equation}
		When $\max\{k_1,k_2\}=k_2$, we can analogously obtain
		\begin{equation}\label{eq:Chaplygin-I2-second-frequency-case}
			\begin{aligned}
				&
				\big\|
				\sum_{\substack{(k_1,k_2)\in\mathcal X_k,\\
				\max\{k_1,k_2\}=k_2}}
				\big\|
				P_k
				\big(
				P_{k_1}\partial_{\alpha\beta}^2\partial_x^b\phi
				P_{k_2}\partial_x^c\cQ^{\alpha\beta}
				\big)
				\big\|_{L^2}
				\big\|_{\ell_k^2}
				\ls
				\|\partial_x\cQ^{\alpha\beta}\|_{H^{|a|-1}}
				\|\partial\phi\|_{B_{\infty,1}^{1+\delta}}.
			\end{aligned}
		\end{equation}
		Note that
		$B_{p,r}^s(\mathbb R^3)\cap L^\infty(\mathbb R^3)$ is an algebra for
		$s>0$ and $p,r\in[1,\infty]$. Then, for $1\le |a|\le N$, one has
		\begin{equation}\label{eq:Chaplygin-Q-Besov}
			\begin{aligned}
				\|\cQ^{\alpha\beta}\|_{B_{\infty,1}^{1+\delta}}
				&\ls
				\|\partial\phi\|_{B_{\infty,1}^{1+\delta}}
				+
				\|\partial\phi\|_{B_{\infty,1}^{1+\delta}}
				\|\partial\phi\|_{L^\infty}
				\\
				&\ls
				\|\partial\phi\|_{B_{\infty,1}^{1+\delta}}
				+
				\|\partial\phi\|_{B_{\infty,1}^{1+\delta}}^2.
			\end{aligned}
		\end{equation}
		Moreover,
		\begin{equation}\label{eq:Chaplygin-Q-high-Sobolev}
			\begin{aligned}
				\|\partial_x\cQ^{\alpha\beta}\|_{H^{|a|-1}}
				\ls
				\|\cQ^{\alpha\beta}\|_{H^{|a|}}
				\ls
				\|\partial\phi\|_{H^{|a|}}
				\left(
				1+
				\|\partial\phi\|_{L^\infty}
				\right).
			\end{aligned}
		\end{equation}
		Hence,
		\begin{equation}\label{eq:Chaplygin-I2-final}
			\|I_2^a\|_{L^2}
			\ls
			\|\partial\phi\|_{H^{|a|}}
			\big(
			\|\partial\phi\|_{B_{\infty,1}^{1+\delta}}
			+
			\|\partial\phi\|_{B_{\infty,1}^{1+\delta}}^2
			\big).
		\end{equation}
	
		
		Collecting \eqref{eq:Chaplygin-basic-energy},
		\eqref{eq:Chaplygin-I1-estimate},
		\eqref{eq:Chaplygin-I2-final}, and using the smallness of
		$\varepsilon_1$, we arrive at
		\begin{equation}\label{eq:Chaplygin-energy-integral}
			\begin{aligned}
				\|\partial\phi(t)\|_{H^N(\mathbb R^3)}^2
				&\ls
				\|\partial\phi(0)\|_{H^N(\mathbb R^3)}^2
				+\int_0^t
				\big(
				\|\partial\phi(\tau)\|_{B_{\infty,1}^{1+\delta}}
				+\|\partial\phi(\tau)\|_{B_{\infty,1}^{1+\delta}}^2
				\big)
				\|\partial\phi(\tau)\|_{H^N(\mathbb R^3)}^2
				d\tau.
			\end{aligned}
		\end{equation}
		
		It follows from \eqref{eq:CVP-bootstrap-pointwise-phi} and $W^{2,\infty}(\R^3) \hookrightarrow B_{\infty,1}^{1+\delta}(\R^3)$ that
		\begin{equation}\label{eq:Chaplygin-time-coefficient}
			\begin{aligned}
				\|\partial\phi(\tau)\|_{B_{\infty,1}^{1+\delta}}
				+\|\partial\phi(\tau)\|_{B_{\infty,1}^{1+\delta}}^2
				\ls
				\varepsilon_1(1+\tau)^{-1}.
			\end{aligned}
		\end{equation}
		Substituting \eqref{eq:Chaplygin-time-coefficient} into
		\eqref{eq:Chaplygin-energy-integral} yields
		\begin{equation}\label{eq:Chaplygin-before-Gronwall}
			\begin{aligned}
				\|\partial\phi(t)\|_{H^N}^2
				&\ls
				\|\partial\phi(0)\|_{H^N}^2
				+
				\int_0^t
				\varepsilon_1(1+\tau)^{-1}
				\|\partial\phi(\tau)\|_{H^N}^2
				d\tau.
			\end{aligned}
		\end{equation}
		By Gronwall's inequality, \eqref{YHC-4} and \eqref{eq:CVP-bootstrap-energy-ordinary},
		\begin{equation}\label{eq:Chaplygin-after-Gronwall}
			\begin{aligned}
				\|\partial\phi(t)\|_{H^N}
				\ls
				\ve + \ve_1^2 \int_{0}^{t} (1+\tau)^{-1+\eta}d\tau
				\ls
				\ve + \ve_1^2 \eta^{-1}(1+t)^\eta.
			\end{aligned}
		\end{equation}
		Therefore, $\ds\sup_{0\le t\le T}(1+t)^{-\eta}\|\partial\phi(t)\|_{H^N(\mathbb R^3)}
			\ls \ve +\ve_1^2$ holds, which yields \eqref{eq:Chaplygin-energy-growth-epsilon1}.
	\end{proof}

\end{document}
